% !TeX root = courtade-kumar.tex
% Included in the three-bit subsection of proof.tex.
For the calculation, let us write $\rho$ and $u$ for the absolute
correlations of $A,B$ and $B,C$, respectively, and let
\begin{equation*}
  f(t)=\kldivb{(1+t)/2}{1/2}.
\end{equation*}
Thus
\begin{equation*}
  \muti{A}{B}=f(\rho),\quad
  \muti{B}{C}=f(u),\quad
  \muti{A}{C}=f(\rho u),\quad
  \muti{A}{B\emid C}=f(\rho)-f(\rho u).
\end{equation*}
Let us also write
\begin{equation*}
  r=\rho^2,\qquad
  \alpha=\frac{(\rho^2-f(\rho))\ln2}{1-\rho^2}.
\end{equation*}
In these parameters, it suffices to prove
\begin{equation}
  \label{eq:three-bit-scalar-target}
  \frac{f(\rho)-f(\rho u)}{2(1+\alpha)}
  +\frac{2\bigl(\rho^2 f(u)-f(\rho u)\bigr)}u
  >\rho^2-f(\rho).
\end{equation}
We will prove this by expanding in $r$ and showing that every
coefficient is positive. For integers $N\ge1$, let us define the divergence tails
\begin{align*}
  b_N&=\sum_{k>N}\frac{1}{2k(2k-1)}
  =\int_0^1\frac{x^{2N}}{1+x}\,\mathrm{d}x,
  &s_N(u)&=\sum_{k>N}\frac{u^{2k}}{2k(2k-1)}
  =\int_0^u\frac{(u-x)x^{2N}}{1-x^2}\,\mathrm{d}x.
\end{align*}
Let us set $b_0=1$ as a separate convention, and define
\begin{align*}
  T_N&=\sum_{i=0}^{N-1}b_i,
  &B_N&=\sum_{i=0}^{N-1}b_i b_{N-i},\\
  G_N(u)&=\sum_{i=0}^{N-1}b_i s_{N-i}(u),
  &C_N(u)&=\sum_{k=1}^N\frac{1-u^{2k}}{2k(2k-1)}.
\end{align*}
The first ingredient compares the divergence tails in two ways.
\begin{lemma}
\label{lem:scalar-tail-order}
For $0<u<1$ and $N\ge1$,
\begin{equation}
  \label{eq:scalar-convolution-order}
  G_N(u)\ge B_N\frac{s_N(u)}{b_N}.
\end{equation}
For $N\ge2$, we also have
\begin{equation}
  \label{eq:scalar-rescaled-order}
  \frac{s_N(u)}{u b_N}\ge\frac{s_2(u^{N/2})}{u^{N/2} b_2}.
\end{equation}
\end{lemma}
\begin{proof}
First, $s_j(u)/b_j$ is the expectation of the decreasing function
$(u-x)_+/(1-x)$ under the probability density proportional to
$x^{2j}/(1+x)$ on $(0,1)$. Increasing $j$ tilts this density by an
increasing function, so $s_j(u)/b_j$ decreases with $j$.
For $1\le j\le N$, therefore,
$s_j(u)\ge b_j s_N(u)/b_N$; summing proves
\eqref{eq:scalar-convolution-order}.
Here is the elementary order principle used in this argument.
If $Z,Z'$ are independent with the same distribution, then
\begin{equation*}
  \operatorname{Cov}(f(Z),g(Z))
  =\frac12\E\sparen{(f(Z)-f(Z'))(g(Z)-g(Z'))}.
\end{equation*}
It is nonnegative when both functions increase and nonpositive when
one increases and the other decreases. Applying this identity to a
density ratio gives the stated comparison of expectations.
For the second comparison, let us fix $a\ge0$ and write
\begin{equation*}
  M_N(a)=\frac{s_N(e^{-a/N})}{e^{-a/N}b_N}.
\end{equation*}
Changing variables in the two tail integrals gives
\begin{equation}
  \label{eq:scalar-tail-kernel}
  M_N(a)=e^{-2a}\E_{\pi_N}
  \sparen{\frac{\sinh(z/N)}{\sinh((z+a)/N)}},
  \qquad
  \pi_N(\mathrm{d}z)\ \propto
  \frac{e^{-2z}}{1+e^{z/N}}\,\mathrm{d}z,
  \quad z>0.
\end{equation}
The quotient inside the expectation increases with $z$ because
$\coth$ decreases. It also increases with $N$, since
$x\coth x$ increases for $x>0$:
its derivative has the sign of $\sinh x\cosh x-x>0$.
Moreover, for $K>N$ the density ratio $\pi_K/\pi_N$ increases with
$z$, as its logarithmic derivative is
\begin{equation*}
  \frac{N^{-1}}{1+e^{-z/N}}
  -\frac{K^{-1}}{1+e^{-z/K}}>0.
\end{equation*}
The order principle and then the pointwise comparison of quotients
show that $M_N(a)$ increases with $N$.
Taking $a=-N\ln u$ gives
$e^{-a/N}=u$ and $e^{-a/2}=u^{N/2}$, proving
\eqref{eq:scalar-rescaled-order}.
All the densities are integrable; at $a=0$ the quotient in
\eqref{eq:scalar-tail-kernel} equals one.
\end{proof}
The second ingredient gives estimates uniform in the coefficient
index $N$. Let us write
\begin{equation*}
  e_N=\frac12\sum_{k=1}^N\frac{1}{2k-1},\qquad
  R_N=2T_N-1,\qquad x=1-u^{2N}.
\end{equation*}

\begin{lemma}
\label{lem:scalar-moments}
For $N\ge2$ and $0\le u\le1$,
\begin{equation*}
  B_N\le b_NR_N,\qquad
  \frac{R_N}{2e_N}<\frac{19}{18},\qquad
  \frac{C_N(u)}{2e_Nb_N}
  \ge 2x+\frac34x^2+\frac18x^3.
\end{equation*}
\end{lemma}
\begin{proof}
The tail integral and finite summation give
\begin{equation*}
  \frac{1}{4N+2}\le b_N\le\frac{1}{4N+1},\qquad
  T_N=1-\ln 2+e_N+Nb_N.
\end{equation*}
Indeed, the integral bounds follow from
$1+x\le2$ and $1+x\ge2\sqrt{x}$.
For positive integers $i,j$, they imply
\begin{equation*}
  b_i^{-1}+b_j^{-1}\ge4(i+j)+2\ge b_{i+j}^{-1},
  \qquad
  b_i b_j\le b_{i+j}(b_i+b_j).
\end{equation*}
Summing over $i+j=N$ and including the $b_0b_N=b_N$ term proves
$B_N\le b_NR_N$.
For the second bound, let us write $t=\sqrt{2/N}\in(0,1]$.
We will use
\begin{equation}
  \label{eq:scalar-harmonic-bound}
  e_N\ge\frac76-\frac t2\ge1-\frac t3.
\end{equation}
The first inequality is equality at $N=2$.
Its right-hand side increases from $N$ to $N+1$ by at most
$e_{N+1}-e_N=1/(4N+2)$, because
\begin{equation*}
  \sqrt{N}\sqrt{N+1}(\sqrt{N}+\sqrt{N+1})
  \ge(2N+1)\sqrt{N}\ge(2N+1)\sqrt2.
\end{equation*}
This proves the first inequality by induction; the second uses
$t\le1$.
We also have
\begin{equation*}
  \frac9{13}<\ln 2<\frac{25}{36}.
\end{equation*}
For the lower bound, the first three terms of
$\ln 2=2\sum_{j\ge0}((2j+1)3^{2j+1})^{-1}$ sum to
$842/1215>9/13$. Replacing $2j+1$ by $3$
for every $j\ge1$ and summing the resulting geometric series gives the upper bound.
Writing $\gamma=\ln 2-1/2>5/26$, we obtain
\begin{align*}
  Nb_N&\le\frac{1}{4+t^2/2}\le\frac14-\frac{t^2}{36},\\
  18Nb_N-e_N
  &\le\frac{10}{3}+\frac{t-t^2}{2}
  \le\frac{83}{24}<\frac{45}{13}<18\gamma.
\end{align*}
Since $R_N=2e_N+2Nb_N-2\gamma$, this is equivalent to
$R_N/(2e_N)<19/18$.
For the third bound, we use the first three nonnegative terms of
the binomial expansion of $1-(1-x)^a$, where $0<a\le1$.
Let us write $s=1/N=t^2/2$ and $e=e_N$. The moment identities
\begin{equation*}
  \sum_{k=1}^N\frac{k}{2(2k-1)}=\frac N4+\frac{e_N}{2},\qquad
  \sum_{k=1}^N\frac{k^2}{2(2k-1)}=\frac{N(N+2)}8+\frac{e_N}{4}
\end{equation*}
follow by dividing the rational summands.
Together with $b_N\le1/(4N+1)$, they yield
\begin{equation}
  \label{eq:scalar-moment-expansion}
  \begin{gathered}
    \frac{C_N(u)}{2e_Nb_N}
    \ge\frac{4+s}{2}\sparen{x+\theta_2x^2+\theta_3x^3},\\
    \theta_2=\frac12-\frac s4-\frac{1}{8e},\qquad
    \theta_3=\frac13-\frac s4+\frac{s^2}{24}
    -\frac{5-2s}{48e}.
  \end{gathered}
\end{equation}
By \eqref{eq:scalar-harmonic-bound},
\begin{equation*}
  e^{-1}\le(1-t/3)^{-1}\le1+t/2.
\end{equation*}
Let
$a_2=(4+s)\theta_2/2-3/4$ and
$a_3=(4+s)\theta_3/2-1/8$.
Substitution gives
\begin{align*}
  a_2&\ge-\frac t8-\frac{5t^2}{32}-\frac{t^3}{64}
  -\frac{t^4}{32}
  \ge-\frac{21t}{64}\ge-\frac t3,\\
  a_3-\nparen{\frac13-\frac t4}
  &\ge\frac{t}{384}
  \sparen{2(1-t)(28-t)+t^2\nparen{(1-t)^2+t^3}}
  \ge0.
\end{align*}
Consequently, the difference between the right-hand side of
\eqref{eq:scalar-moment-expansion} and
$2x+3x^2/4+x^3/8$ is bounded below by
\begin{align*}
  x\sparen{\frac{t^2}{4}-\frac{tx}{3}
  +\nparen{\frac13-\frac t4}x^2}
  &=\frac{x}{12}
  \sparen{3(t-x)^2+2tx(1-x)+x^2(1-t)}\\
  &\ge0.
\end{align*}
This proves the last bound.
\end{proof}
The preceding estimates will reduce every coefficient of index at
least two to the function
\begin{equation*}
  \mathcal J(w)=\frac{55}{72}-\frac{31w^4}{8}
  +\frac{9w^8}{8}-\frac{w^{12}}8
  +38\frac{s_2(w)}w,\qquad 0<w<1.
\end{equation*}

\begin{lemma}
\label{lem:scalar-fixed-function}
For $0<w<1$, we have $\mathcal J(w)>0$.
\end{lemma}
\begin{proof}
Let us write $g(w)=s_2(w)/w$. Differentiation gives
\begin{align*}
  g'(w)&=-\frac{\ln(1-w^2)}{2w^2}-\frac12-\frac{w^2}{4}
  =\sum_{k\ge3}\frac{w^{2k-2}}{2k},\\
  g''(w)&=\frac{1}{w(1-w^2)}
  +\frac{\ln(1-w^2)}{w^3}-\frac w2.
\end{align*}
Thus
\begin{equation*}
  \frac{\mathcal J'(w)}{w^3}
  =-\frac{31}{2}+9w^4-\frac32w^8
  +19\sum_{k\ge3}\frac{w^{2k-5}}{k}
\end{equation*}
is strictly increasing on $(0,1)$.
Also $g'''(w)>0$, and for $w\ge4/5$,
\begin{equation*}
  \nparen{-\frac{31w^4}{8}+\frac{9w^8}{8}
  -\frac{w^{12}}8}'''
  =w(-93+378w^4-165w^8)>0.
\end{equation*}
Indeed, the quadratic $-93+378z-165z^2$ increases on $[0,1]$
and is positive at $z=2/5<(4/5)^4$.
Hence $\mathcal J'''(w)>0$ for $w\ge4/5$.
We need four finite estimates:
\begin{equation}
  \label{eq:scalar-fixed-checks}
  \mathcal J'(4/5)<0,\qquad
  \mathcal J''(4/5)>25,\qquad
  \mathcal J(5/6)>\frac1{256},\qquad
  |\mathcal J'(5/6)|<\frac15.
\end{equation}
For completeness, the following rational bounds certify them.
For $z>1$, let us write $v=(z-1)/(z+1)$ and define
\begin{equation*}
  P(z)=2\sum_{j=0}^5\frac{v^{2j+1}}{2j+1},\qquad
  E(z)=\frac{2v^{13}}{13(1-v^2)}.
\end{equation*}
The logarithm series gives
$P(z)<\ln z<P(z)+E(z)$.
Only $z\in\{2,3/2,5/4,11/8\}$ are needed, using
\begin{align*}
  \ln(9/25)&=2\ln(3/2)-2\ln(5/4)-2\ln 2,\\
  \ln(11/36)&=\ln(11/8)-2\ln(3/2)-\ln 2,\\
  (\ln 2)\kldivb{11/12}{1/2}&=\frac34\ln 2+\frac{11}{12}\ln(11/8)-\ln(3/2).
\end{align*}
Substituting the appropriate lower or upper bound for each logarithm
in the displayed formulas for $g',g''$, and
\begin{equation*}
  g(w)=\frac{(\ln 2)\kldivb{(1+w)/2}{1/2}-w^2/2-w^4/12}{w},
\end{equation*}
gives
\begin{equation*}
  \begin{array}{c|c|c}
    \text{quantity}&\text{strict lower bound}&\text{strict upper bound}\\ \hline
    \mathcal J'(4/5)&-1&-9/10\\
    \mathcal J''(4/5)&25&26\\
    \mathcal J(5/6)&9/2000&1/200\\
    \mathcal J'(5/6)&9/50&19/100
  \end{array}
\end{equation*}
These are finite rational comparisons and imply
\eqref{eq:scalar-fixed-checks}.
Since $\mathcal J'''$ is positive on $[4/5,1)$, we have
$\mathcal J''>25$ there. Taylor's formula at $5/6$ therefore gives
\begin{align*}
  \mathcal J(w)
  &\ge\mathcal J(5/6)+\mathcal J'(5/6)(w-5/6)
  +\frac{25}{2}(w-5/6)^2\\
  &>\frac1{256}-\frac1{1250}>0
  \qquad(4/5\le w<1).
\end{align*}
On $(0,4/5]$, the increasing ratio $\mathcal J'(w)/w^3$
and the inequality $\mathcal J'(4/5)<0$ show that
$\mathcal J$ decreases. Its value there is consequently at least
$\mathcal J(4/5)>0$.
\end{proof}
We now apply these estimates to the coefficients
\begin{equation*}
  c_N(u)=C_N(u)+\frac{4G_N(u)}u-2B_N,
  \qquad N\ge1,\quad 0<u\le1.
\end{equation*}

\begin{lemma}
\label{lem:scalar-coefficients}
For every integer $N\ge1$ and $0<u\le1$, we have $c_N(u)>0$.
\end{lemma}
\begin{proof}
For $N\ge2$ and $0<u<1$, let us write
\begin{equation*}
  k=\frac{4s_N(u)}{u b_N}-2,\qquad
  w=u^{N/2},\qquad x=1-w^4.
\end{equation*}
By \autoref{lem:scalar-tail-order},
$c_N(u)\ge C_N(u)+B_Nk$.
If $k\ge0$, this is positive because $C_N(u)>0$.
If $k<0$, \autoref{lem:scalar-moments} gives
\begin{align*}
  \frac{c_N(u)}{2e_Nb_N}
  &\ge\frac{C_N(u)}{2e_Nb_N}+\frac{R_N}{2e_N}k\\
  &\ge2x+\frac34x^2+\frac18x^3+\frac{19}{18}k.
\end{align*}
Here the upper bounds on $B_N$ and $R_N/(2e_N)$ give lower bounds
after multiplication by the negative number $k$.
Since $b_2=\ln 2-7/12<1/9$, the rescaled tail comparison gives
\begin{equation*}
  k\ge36\frac{s_2(w)}w-2.
\end{equation*}
Substituting this and $x=1-w^4$ into the preceding lower bound
gives exactly $\mathcal J(w)>0$.
For $u=1$, we instead have $s_j(1)=b_j$ and $C_N(1)=0$,
so $c_N(1)=2B_N>0$.
It remains to treat $N=1$.
Writing $\gamma=\ln 2-1/2<7/36$ and retaining the first two terms
of $s_1(u)$ gives
\begin{align*}
  c_1(u)&=\frac{1-u^2}{2}+\frac{4s_1(u)}u-2\gamma\\
  &\ge Z(u):=\frac19-\frac{u^2}{2}
  +\frac{u^3}{3}+\frac{2u^5}{15}.
\end{align*}
To see that $Z>0$, note that
$Z'(u)/u=-1+u+2u^3/3$ increases and is negative at $u=2/3$.
Also $Z''(u)=-1+2u+8u^3/3>1$ on $[2/3,1]$.
At $a=3/4$,
\begin{equation*}
  Z(a)=\frac{49}{23040}>\frac1{512},\qquad
  Z'(a)=\frac3{128}<\frac1{32}.
\end{equation*}
Taylor's formula on $[2/3,1]$ yields
\begin{equation*}
  Z(u)\ge Z(a)+Z'(a)(u-a)+\frac12(u-a)^2
  >\frac1{512}-\frac1{2048}>0.
\end{equation*}
On $(0,2/3]$, the monotonicity of $Z'(u)/u$ makes $Z$
decreasing, so $Z(u)\ge Z(2/3)>0$ there as well.
\end{proof}

\begingroup
\newtheorem*{restatedthreebit}{Lemma~\ref{lem:three-bit}}
\begin{restatedthreebit}[restated]
\begin{align*}
  &\frac{\muti{A}{B\emid C}}
  {2\left(1+\ln2\,
  \frac{\tsmuti{A}{B}-\muti{A}{B}}{\tsent{A\emid B}}\right)}
  +
  \frac{2\bigl[\tsmuti{A}{B}\muti{B}{C}-\muti{A}{C}\bigr]}
  {\sqrt{\tsmuti{B}{C}}}
  \ge\tsmuti{A}{B}-\muti{A}{B}.
\end{align*}
\end{restatedthreebit}
\endgroup
\begin{proof}
The divergence-tail definitions give the absolutely convergent expansions
\begin{align*}
  \alpha&=\sum_{N\ge1}b_Nr^N,\\
  \frac{\ln2}{1-\rho^2}
  \bigl[\rho^2 f(u)-f(\rho u)\bigr]
  &=\sum_{N\ge1}s_N(u)r^N,\\
  \frac{\ln2}{1-\rho^2}
  \bigl[f(\rho)-f(\rho u)\bigr]
  &=\sum_{N\ge1}C_N(u)r^N.
\end{align*}
In particular, $\alpha>0$.
Multiplying these series, with the convention $b_0=1$, gives
\begin{equation*}
  \begin{aligned}
    &\frac{\ln2}{1-\rho^2}\Biggl[
    f(\rho)-f(\rho u)
    +\frac{4(1+\alpha)\bigl(\rho^2 f(u)-f(\rho u)\bigr)}u\\
    &\hspace{35mm}
    -2\bigl(\rho^2-f(\rho)\bigr)(1+\alpha)\Biggr]
    =\sum_{N\ge1}c_N(u)r^N.
  \end{aligned}
\end{equation*}
Every coefficient on the right is positive by
\autoref{lem:scalar-coefficients}. Since $0<r<1$, the expression
in brackets is positive. Dividing it by $2(1+\alpha)$ proves
\eqref{eq:three-bit-scalar-target}.
\end{proof}

